\documentclass[11pt]{article}
\usepackage[margin=1in]{geometry}
\usepackage{amsmath,amssymb,amsthm}
\usepackage{hyperref}
\usepackage{xurl}

\newtheorem{theorem}{Theorem}
\newtheorem{lemma}[theorem]{Lemma}
\theoremstyle{definition}
\newtheorem{definition}[theorem]{Definition}
\theoremstyle{remark}
\newtheorem{remark}[theorem]{Remark}

\let\oldus\_
\renewcommand{\_}{\oldus\allowbreak}

\title{A proof of Conjecture 00000004930\\[2pt]
\large Two MDPs with equal discounted values for every discount factor but different average
rewards}
\author{Jackmeson1}
\date{}

\begin{document}
\sloppy
\maketitle

\section{The conjecture}

Conjecture 00000004930 reads (English version; the Chinese version says the same, with
``all discount-factor values identical''):

\begin{quote}
\emph{Definition:} Discounted value and average reward are two layers of optimality measures.
\emph{Conjecture:} There exist two MDPs with identical values for all discount factors but
different average rewards, and the separation is realized by an explicit construction
concentrating the discount weights on transients.
\end{quote}

It is an existence statement. We give two explicit Markov decision processes with countable state
spaces whose optimal discounted values agree for every discount factor, while under every policy
the long-run average reward is $1$ in the first and $0$ in the second. The second MDP collects all
of its reward on transient states. Everything is proved in Lean 4 with Mathlib.

\section{Definitions}

\begin{definition}[MDP, policies]
An MDP $M$ (Lean: \texttt{MDP}) consists of a state type $S$, an action type $\mathcal A$, an initial
state $s_0$, transition laws $P(s,a)\in\mathrm{PMF}(S)$ and rewards $r(s,a)\in\mathbb R_{\ge0}$. A
history at time $t$ is $((s_0,a_0),\dots,(s_{t-1},a_{t-1});s_t)$ (Lean: \texttt{Hist}). A policy is
history-dependent and randomized: $\pi(h)\in\mathrm{PMF}(\mathcal A)$ (Lean: \texttt{Policy}). The law
of the history at time $t$ is defined recursively by $s_0$ at $t=0$ and, from a history $h$ with
current state $s$, by drawing $a\sim\pi(h)$, $s'\sim P(s,a)$ and appending $(s,a)$ (Lean:
\texttt{law}).
\end{definition}

\begin{definition}[Values and average reward]
The expected reward at time $t$ is $\mathbb E^\pi[r(s_t,a_t)]$ (Lean: \texttt{expReward}). For a
discount factor $\beta$, the discounted value of $\pi$ is $V_\beta(\pi)=\sum_{t\ge0}\beta^t\,
\mathbb E^\pi[r(s_t,a_t)]$ (Lean: \texttt{discValue}), and the (optimal) value of $M$ is
$V^*_\beta=\sup_\pi V_\beta(\pi)$ (Lean: \texttt{optValue}). The $T$-step average is
$\frac1T\sum_{t<T}\mathbb E^\pi[r(s_t,a_t)]$ (Lean: \texttt{avgReward}); the average reward of $\pi$
is its limit (or liminf, or limsup) as $T\to\infty$, and the optimal average reward is the supremum
over policies. Lean computes the $T$-step average from the real numbers
\texttt{(expReward M $\pi$ t).toReal}; since \texttt{toReal} sends $\infty$ to $0$, this agrees with
the usual definition whenever the expected rewards are finite, e.g.\ for bounded rewards. Both
witnesses have rewards in $\{0,1\}$, so their expected rewards are at most $1$
(\texttt{expReward\_le\_one}) and the formal and usual averages coincide.

Values are taken in $[0,\infty]$. The Lean statements let $\beta$ range over all of $[0,\infty]$ with
the arithmetic of $[0,\infty]$, in which $1-\beta$ is truncated subtraction; so the Lean value
$(1-\beta)^{-1}$ means $1/(1-\beta)$ for a discount factor $\beta\in[0,1)$ and $\infty$ for
$\beta\ge1$ (where the series $\sum_t\beta^t$ indeed diverges). Below we write $(1-\beta)^{-1}$ in
this sense; for every discount factor $\beta\in[0,1)$ it is the ordinary number $1/(1-\beta)$.
\end{definition}

\section{The two MDPs}

\textbf{MDP $A$} (Lean: \texttt{A}): one state, one action, reward $1$.

\textbf{MDP $B$} (Lean: \texttt{B}): states $\mathit{start}$, $\mathit{count}(k)$ for $k\in\mathbb N$, and
$\mathit{sink}$; actions $n\in\mathbb N$; deterministic transitions (Lean: \texttt{next})
\[ \mathit{start}\xrightarrow{n}\mathit{count}(n),\quad \mathit{count}(k+1)\to\mathit{count}(k),\quad
\mathit{count}(0)\to\mathit{sink},\quad \mathit{sink}\to\mathit{sink}, \]
where the action matters only at $\mathit{start}$; reward $1$ at every state except $\mathit{sink}$,
where it is $0$. Playing $n$ first yields reward $1$ at times $0,1,\dots,n+1$ and $0$ afterwards.
All rewarded states are transient: each is visited at most once, and the discount weights
$\beta^t$ that make up $V_\beta$ are carried by this transient prefix.

\section{Results}

\begin{lemma}[\texttt{optValue\_A}, \texttt{avgReward\_A}]
In $A$ every policy has expected reward $1$ at every time. Hence $V^*_\beta(A)=(1-\beta)^{-1}$ for
every $\beta$, and every policy has average reward $1$.
\end{lemma}

\begin{lemma}[\texttt{optValue\_B}]
$V^*_\beta(B)=(1-\beta)^{-1}$ for every $\beta$.
\end{lemma}

\begin{proof}
Rewards are at most $1$, so the expected reward is at most $1$ (\texttt{expReward\_le\_one}) and
$V_\beta(\pi)\le\sum_t\beta^t$. Conversely the deterministic policy that always plays $n$
(\texttt{detPol}) has expected reward $1$ at times $t\le n+1$ (\texttt{expReward\_detPol}), so
$V^*_\beta(B)\ge\sum_{t<n}\beta^t$ for every $n$, whose supremum is $\sum_t\beta^t=(1-\beta)^{-1}$.
\end{proof}

\begin{lemma}[\texttt{law\_state}]
Let $\pi$ be any policy for $B$ and $q=\pi(\mathit{start})$ the law of the first action. For every
$t\ge1$ the law of the state $s_t$ is the image of $q$ under $n\mapsto \mathrm{st}(n,t)$, where
$\mathrm{st}(n,\cdot)$ is the state sequence obtained by playing $n$.
\end{lemma}

\begin{proof}
Induction on $t$. For $t=1$ this is the first transition. If the claim holds at $t\ge1$, then every
history in the support of the time-$t$ law has state $\ne\mathit{start}$; at such states the
transition does not depend on the action, so the randomized, history-dependent choice of $a_t$ is
irrelevant and the state law is pushed forward by the fixed step map.
\end{proof}

\begin{lemma}[\texttt{expReward\_B\_le}, \texttt{expReward\_B\_tendsto}, \texttt{avgReward\_B}]
Under every policy for $B$: $\mathbb E^\pi[r_{t+1}]\le\sum_{k}q(k+t)$, so the expected reward
tends to $0$, and by Cesaro's theorem the average reward is $0$.
\end{lemma}

\begin{theorem}[\texttt{conjecture4930}]
For every $\beta$: $V^*_\beta(A)=V^*_\beta(B)=(1-\beta)^{-1}$. Under every policy the average reward
of $A$ tends to $1$ and that of $B$ tends to $0$; in $B$ the expected reward at time $t$ tends to
$0$ under every policy. Hence $A$ and $B$ have identical values for all discount factors and
different (optimal) average rewards, $1\neq0$.
\end{theorem}

\section{Remarks on the reading}

\begin{itemize}
\item \emph{Values} are optimal values at the initial state, and \emph{average rewards} are long-run
average rewards; since the limit exists for every policy, the lim, liminf and limsup versions and
their suprema over policies all give $1$ for $A$ and $0$ for $B$.
\item The separation needs a countably infinite state space. For finite MDPs it is impossible,
because $g^*=\lim_{\beta\uparrow1}(1-\beta)V^*_\beta$ (Blackwell); for a single fixed policy it is
impossible as well, because $\beta\mapsto V_\beta$ is the generating function of the expected
reward sequence, which it therefore determines. The conjecture does not restrict to finite MDPs,
and its phrase ``concentrating the discount weights on transients'' describes exactly this
construction: all reward of $B$ lies on transient states, which the discounted criterion sees and
the average criterion ignores.
\item Rewards are taken nonnegative (both witnesses have rewards in $\{0,1\}$); policies are fully
general (history-dependent and randomized).
\end{itemize}

\section{Lean formalization and axiom audit}

The file \texttt{lean/Conjecture4930/Basic.lean} (300 lines, Lean 4.33.1, Mathlib v4.33.1) contains
the definitions above and the theorems \texttt{optValue\_A}, \texttt{avgReward\_A},
\texttt{optValue\_B}, \texttt{law\_state}, \texttt{expReward\_B\_le}, \texttt{expReward\_B\_tendsto},
\texttt{avgReward\_B} and the main theorem \texttt{Conjecture4930.conjecture4930}. It builds with
\texttt{lake build}. \texttt{\#print axioms conjecture4930} (and \texttt{law\_state}) reports exactly
\texttt{[propext, Classical.choice, Quot.sound]}. There is no \texttt{sorry}, no
\texttt{native\_decide} and no added axiom.

\end{document}
